\documentclass[11pt]{article}

% -----------------------------------------------------------------------
% Packages
% -----------------------------------------------------------------------
\usepackage{amsmath,amssymb,amsthm}
\usepackage{geometry}
\usepackage{hyperref}
\usepackage{graphicx}
\usepackage{booktabs}
\usepackage{xcolor}
\usepackage{authblk}

\geometry{margin=1.2in}

% -----------------------------------------------------------------------
% Theorem environments
% -----------------------------------------------------------------------
\newtheorem{theorem}{Theorem}[section]
\newtheorem{proposition}[theorem]{Proposition}
\newtheorem{lemma}[theorem]{Lemma}
\newtheorem{corollary}[theorem]{Corollary}
\newtheorem{remark}[theorem]{Remark}
\newtheorem{definition}[theorem]{Definition}

% -----------------------------------------------------------------------
% Macros
% -----------------------------------------------------------------------
\newcommand{\keff}{k_{\mathrm{eff}}}
\newcommand{\lc}{\lambda_c}
\newcommand{\rhon}{\rho_n}
\newcommand{\Sigmar}{\Sigma_r}
\newcommand{\nuSigmaf}{\nu\Sigma_f}
\newcommand{\R}{\mathbb{R}}
\newcommand{\N}{\mathbb{N}}
\newcommand{\Z}{\mathbb{Z}}

% -----------------------------------------------------------------------
% Title block
% -----------------------------------------------------------------------
\title{\textbf{Criticality Thresholds in One-Dimensional Multiplying Media
       with $n$-Bonacci Aperiodic Modulation}\\[6pt]
       {\large Spectral Gap Control of $k_{\mathrm{eff}}$ in
        Substitution-Sequence Diffusion Operators}}

\author[1]{Pablo Nogueira Grossi}
\affil[1]{G6 LLC, Newark, New Jersey, USA\\
          \texttt{pablogrossi@hotmail.com}\\
          ORCID: 0009-0000-6496-2186}

\date{2026}

% -----------------------------------------------------------------------
\begin{document}
% -----------------------------------------------------------------------

\maketitle

\begin{abstract}
We study the one-group neutron diffusion equation on a uniform
one-dimensional slab whose material coefficients---diffusion constant,
removal cross-section, and fission production rate---vary site-by-site
according to the $n$-bonacci substitution sequence for
$n = 2, 3, 4, 5$.
The criticality condition $\keff = 1$ defines a critical fission strength
$\lc(n)$, which we compute by solving the associated generalized
eigenvalue problem via finite differences.
We find that $\lc(n)$ is governed not by the $n$-bonacci constant $\rhon$
alone, but by the \emph{spectral gap} of the substitution transfer matrix,
$\Delta_n = \rhon - |\rho_n^{(2)}|$, where $\rho_n^{(2)}$ is the
subdominant root of the characteristic polynomial
$x^n = x^{n-1} + \cdots + 1$.
Numerically, $\lc(n) \approx 0.958\,\Delta_n + 0.107$ with correlation
$r = 0.989$ across $n = 2, \ldots, 5$.
For $n \geq 4$, $\lc(n)$ converges to $7/6$ exactly;
the Tribonacci case $n=3$ saturates at the distinct value
$\lc(3) \approx 37/32$, while $\lc(2) \approx 1.064$ for Fibonacci.
These results extend and complement a companion study of nonlinear
(DNLS) dynamics on Fibonacci and Tribonacci chains~\cite{Grossi2026DNLS},
and provide a concrete spectral-theoretic framework for aperiodic
multiplying media.
\end{abstract}

\tableofcontents
\bigskip

% ======================================================================
\section{Introduction}
\label{sec:intro}
% ======================================================================

The interplay between aperiodic order and spectral properties of
differential operators has been a central theme in mathematical physics
since the discovery of quasicrystals~\cite{Shechtman1984} and the
subsequent rigorous treatment of almost-periodic Schr\"{o}dinger
operators~\cite{Bellissard1992,Damanik2017}.
The Fibonacci chain---generated by the substitution
$A \mapsto AB$, $B \mapsto A$---is the canonical one-dimensional model:
its tight-binding Hamiltonian has a Cantor-set spectrum of measure zero,
multifractal eigenstates, and spectral gaps labelled by the gap-labelling
theorem~\cite{Bellissard1992}.
The Tribonacci chain, generated by $A \mapsto AB$, $B \mapsto AC$,
$C \mapsto A$, extends this framework to three letter alphabets and
introduces richer hierarchical structure governed by the tribonacci
constant $\rho_3 \approx 1.839$~\cite{Krebbekx2023}.

Transport theory in heterogeneous media provides a physically motivated
arena for the same mathematical structures.
The neutron diffusion equation in a multiplying medium,
\begin{equation}
  -\nabla\cdot\bigl(D(\mathbf{r})\,\nabla\phi\bigr)
  + \Sigmar(\mathbf{r})\,\phi
  \;=\;
  \frac{1}{\keff}\,\nuSigmaf(\mathbf{r})\,\phi,
  \label{eq:diffusion-strong}
\end{equation}
defines $\keff$ as the dominant eigenvalue of the ratio of fission
production to neutron loss.
When the material coefficients $D$, $\Sigmar$, $\nuSigmaf$ are spatially
uniform, $\keff$ is determined by simple geometric buckling.
When they vary periodically, Bloch theory applies and band-gap
phenomena are well understood~\cite{Squires1978}.
When they vary \emph{aperiodically}---as in a substitution-sequence
medium---neither framework is directly applicable, and the spectral theory
of the associated operator must be studied on its own terms.

\subsection{Motivation and relation to prior work}

This paper is motivated by two converging lines of inquiry.

The first is a companion study~\cite{Grossi2026DNLS} in which the
discrete nonlinear Schr\"{o}dinger (DNLS) equation is evolved on
Fibonacci and Tribonacci hopping chains.
There, the mid-gap eigenstates are multifractal---a hallmark of
critical (neither localized nor extended) behavior---and the
inverse participation ratio (IPR) responds differentially to
nonlinearity: Fibonacci and Tribonacci chains exhibit distinct
robustness thresholds as the coupling $\lambda$ increases.
The present work asks the same structural question in a transport setting:
\emph{how does the substitution rule governing spatial heterogeneity
control the criticality threshold of a multiplying medium?}

The second motivation is the gap-labelling program for one-dimensional
operators with substitution-sequence potentials.
The spectral gap between the dominant and subdominant eigenvalues of the
transfer matrix---what we call $\Delta_n = \rhon - |\rho_n^{(2)}|$---governs
the rate at which finite-generation approximants converge to the
thermodynamic limit.
We find, numerically, that $\Delta_n$ also directly controls $\lc(n)$,
suggesting a precise connection between the transfer-matrix spectral gap
and the physical criticality threshold.
This connection is the central result of the paper.

\subsection{Main results}

We establish the following, numerically for $n = 2, 3, 4, 5$ and
analytically for the linear relationship in the large-$N$ limit:

\begin{enumerate}
  \item \textbf{Spectral gap governs criticality.}
    The critical fission strength $\lc(n)$ satisfies
    \begin{equation}
      \lc(n) \;\approx\; \alpha\,\Delta_n + \beta,
      \qquad
      \alpha \approx 0.958,\quad \beta \approx 0.107,
      \label{eq:linear-fit}
    \end{equation}
    with correlation $r = 0.989$ across the $n$-bonacci family $n=2,\ldots,5$.

  \item \textbf{Saturation and exact values.}
    For $n \geq 4$, $\lc(n) \to 7/6$ exactly (verified to machine
    precision for $n=5$, converging from below for $n=4$).
    The Tribonacci case $n=3$ saturates at the distinct value
    $\lc(3) \approx 37/32 = 1.15625$, stable across generations $g = 7$--$9$.
    The Fibonacci case $n=2$ gives $\lc(2) \approx 1.064$, the lowest threshold,
    reflecting the minimal absorber content of the two-symbol alphabet.
    The existence of exact rational limits for $n \geq 3$ suggests
    an analytic explanation via the transfer-matrix fixed-point structure.

  \item \textbf{Fibonacci as outlier.}
    The two-letter Fibonacci substitution ($n=2$) occupies a qualitatively
    distinct regime from the $n \geq 3$ family.
    Its spectral gap $\Delta_2 = \rho_2 - |\rho_2^{(2)}| = 1$ is the only
    integer value among $n = 2,\ldots,5$, and its critical threshold
    $\lc(2) \approx 1.064$ is the lowest, reflecting the minimal absorber
    content of the two-symbol alphabet.

  \item \textbf{Generation convergence.}
    $\keff$ converges to its thermodynamic-limit value with generation $g$
    at a rate consistent with the spectral gap of the substitution matrix.
    Tribonacci ($n=3$) converges more slowly than Fibonacci ($n=2$),
    consistent with its smaller $\Delta_n$.
\end{enumerate}

\subsection{Organization}

Section~\ref{sec:model} introduces the $n$-bonacci substitution sequences,
the one-group diffusion model, and the finite-difference discretization.
Section~\ref{sec:spectral} reviews the transfer-matrix spectral theory
that motivates the $\lc$--$\Delta_n$ relationship.
Section~\ref{sec:numerics} presents the numerical results: the $\lc(n)$
table, the spectral gap correlation, the saturation phenomenon, and
generation convergence.
Section~\ref{sec:flux} examines the fundamental flux modes and their
hierarchical spatial structure.
Section~\ref{sec:discussion} discusses connections to the DNLS companion
paper, the gap-labelling theorem, and open questions.

% ======================================================================
\section{Model}
\label{sec:model}
% ======================================================================

% ----------------------------------------------------------------------
\subsection{The $n$-bonacci substitution sequences}
\label{sec:nbonacci}
% ----------------------------------------------------------------------

For each integer $n \geq 2$, the $n$-bonacci substitution acts on an
alphabet $\mathcal{A}_n = \{A_0, A_1, \ldots, A_{n-1}\}$ by the rules
\begin{equation}
  A_0 \;\mapsto\; A_0 A_1,
  \qquad
  A_k \;\mapsto\; A_0 A_{k+1} \quad (1 \leq k \leq n-2),
  \qquad
  A_{n-1} \;\mapsto\; A_0.
  \label{eq:substitution}
\end{equation}
Beginning from the seed word $A_0$, iterating the substitution $g$ times
produces a word $\mathbf{w}^{(n,g)}$ of length $|\mathbf{w}^{(n,g)}|$
that grows as $\rhon^g$ where $\rhon$ is the unique real root greater than
one of the characteristic polynomial
\begin{equation}
  p_n(x) \;=\; x^n - x^{n-1} - x^{n-2} - \cdots - x - 1.
  \label{eq:charpoly}
\end{equation}

\begin{definition}[$n$-bonacci constant]
  The \emph{$n$-bonacci constant} $\rhon$ is the unique positive real root
  of $p_n(x) = 0$ with $\rhon > 1$.
\end{definition}

The cases $n = 2$ and $n = 3$ recover the familiar Fibonacci
($\rho_2 = (1+\sqrt{5})/2 \approx 1.618$) and Tribonacci
($\rho_3 \approx 1.839$) substitutions respectively.
Table~\ref{tab:nbonacci} collects the constants and chain lengths
used in this paper.

\begin{table}[h]
\centering
\caption{$n$-bonacci constants, subdominant root moduli, spectral gaps,
         and chain lengths at generation $g$ used in numerical experiments.}
\label{tab:nbonacci}
\begin{tabular}{ccccccc}
\toprule
$n$ & $\rhon$ & $|\rho_n^{(2)}|$ & $\Delta_n$ & $g$ & $N = |\mathbf{w}^{(n,g)}|$ \\
\midrule
2 & 1.61803 & 0.61803 & 1.00000 & 10 & 144 \\
3 & 1.83929 & 0.73735 & 1.10193 &  8 & 149 \\
4 & 1.92756 & 0.81828 & 1.10929 &  8 & 208 \\
5 & 1.96595 & 0.87105 & 1.09490 &  8 & 236 \\
\bottomrule
\end{tabular}
\end{table}

% ----------------------------------------------------------------------
\subsection{The one-group diffusion model}
\label{sec:diffusion}
% ----------------------------------------------------------------------

We consider a one-dimensional slab of length $L = N\,h$ (where $h$ is the
mesh spacing, set to $h = 1$ throughout) occupied by $N$ sites.
Each site $i \in \{1, \ldots, N\}$ carries material parameters determined
by the symbol $w_i \in \mathcal{A}_n$ at position $i$ in the substitution
word $\mathbf{w}^{(n,g)}$:

\begin{equation}
  (D_i,\,\Sigmar^{(i)},\,\nuSigmaf^{(i)})
  \;=\;
  \begin{cases}
    (1.0,\;0.5,\;\lambda) & \text{if } w_i = A_0 \text{ (fissile)}, \\
    (1.0,\;2.0,\;0)       & \text{if } w_i \in \{A_1,\ldots,A_{n-1}\} \text{ (absorber)}.
  \end{cases}
  \label{eq:parammap}
\end{equation}

The parameter $\lambda \geq 0$ is the fission production strength and
is the sole free parameter of the model.
All other values are dimensionless and fixed; they are chosen so that
(i) fissile and absorber sites differ sharply in both removal cross-section
($\Sigmar^{(0)} = 0.5$ vs.\ $\Sigmar^{(k)} = 2.0$) and fission content,
creating genuine spectral competition between production and loss, and
(ii) the system crosses criticality ($\keff = 1$) at a value of $\lambda$
of order unity, making finite-size effects observable across reasonable
generation numbers.

\begin{remark}
  The model can be interpreted as a one-speed diffusion approximation in a
  heterogeneous multiplying medium with aperiodically modulated coefficients,
  in the spirit of quasicrystal-inspired transport problems.
  The substitution word encodes a deterministic aperiodic disorder rather
  than a random or periodic one; this is the key distinction from classical
  homogenization theory.
\end{remark}

The one-group diffusion equation~\eqref{eq:diffusion-strong} on this
domain, with vacuum boundary conditions $\phi(0) = \phi(L) = 0$,
becomes the eigenvalue problem
\begin{equation}
  L\,\boldsymbol{\phi} \;=\; \frac{1}{\keff}\,F\,\boldsymbol{\phi},
  \label{eq:evp}
\end{equation}
where $\boldsymbol{\phi} = (\phi_1,\ldots,\phi_N)^T$ is the discretized
flux vector, $L$ is the loss (diffusion plus removal) matrix, and $F$ is
the fission matrix defined below.

% ----------------------------------------------------------------------
\subsection{Finite-difference discretization}
\label{sec:fd}
% ----------------------------------------------------------------------

We discretize equation~\eqref{eq:diffusion-strong} by standard
cell-centered finite differences on the uniform mesh $x_i = (i-1/2)h$,
$i = 1,\ldots,N$.
At interior cell interfaces $x_{i+1/2}$, the diffusion coefficient is
approximated by the harmonic mean
\begin{equation}
  D_{i+1/2}
  \;=\;
  \frac{2\,D_i\,D_{i+1}}{D_i + D_{i+1}},
  \label{eq:harmonic}
\end{equation}
which is the standard choice for discontinuous diffusion coefficients
as it ensures continuity of the neutron current $J = -D\,d\phi/dx$
across material interfaces~\cite{Duderstadt1976}.
At the boundaries, vacuum conditions are implemented by setting
$D_{1/2} = D_1$ and $D_{N+1/2} = D_N$.

The loss matrix $L \in \R^{N \times N}$ is symmetric tridiagonal:
\begin{align}
  L_{ii}   &= D_{i-1/2} + D_{i+1/2} + \Sigmar^{(i)}, \label{eq:L-diag}\\
  L_{i,i\pm 1} &= -D_{i\pm 1/2}.                      \label{eq:L-offdiag}
\end{align}
The fission matrix $F = \mathrm{diag}(\nuSigmaf^{(1)},\ldots,\nuSigmaf^{(N)})$
is diagonal and non-negative.
Since the only non-zero entries of $F$ occur at fissile sites ($w_i = A_0$),
$F$ is in general rank-deficient, with $\mathrm{rank}(F)$ equal to the
number of $A_0$ symbols in $\mathbf{w}^{(n,g)}$.

The eigenvalue problem~\eqref{eq:evp} is equivalent to
\begin{equation}
  F\,\boldsymbol{\phi} \;=\; \keff\,L\,\boldsymbol{\phi},
  \label{eq:gevp}
\end{equation}
which is a generalized eigenvalue problem with $L$ symmetric positive
definite and $F$ symmetric positive semidefinite.
By the Perron--Frobenius theorem applied to $L^{-1}F$, the dominant
eigenvalue $\keff$ is real, positive, and simple, with a strictly
positive eigenvector (the fundamental flux mode)~\cite{Stacey2007}.

We compute $\keff$ and the fundamental mode using \texttt{scipy.linalg.eig}
applied to the dense matrices, which is feasible for $N \leq 300$.
For reproducibility, all code is available at~\cite{GrossiRepo2026}.

% ----------------------------------------------------------------------
\subsection{The critical fission strength $\lc(n)$}
\label{sec:lambdac}
% ----------------------------------------------------------------------

For fixed $n$ and generation $g$, $\keff(\lambda)$ is a strictly increasing,
continuous function of $\lambda$ (since increasing the fission production at
fissile sites monotonically increases the dominant eigenvalue of $L^{-1}F$).
We define the \emph{critical fission strength} as the unique root of
\begin{equation}
  \keff(\lambda)\;=\;1,
  \label{eq:criticality}
\end{equation}
which is computed by Brent's method~\cite{Brent1973} bracketed on
$[\lambda_{\min}, \lambda_{\max}] = [0.3, 6.0]$.

The quantity $\lc(n,g)$ depends on $g$ through the finite-size geometry of
$\mathbf{w}^{(n,g)}$, and converges as $g \to \infty$ to a
thermodynamic-limit value $\lc(n) := \lim_{g\to\infty}\lc(n,g)$.
We study this convergence in Section~\ref{sec:numerics} and use the
largest tractable generation as a proxy for the thermodynamic limit.

% ======================================================================
% Bibliography placeholder
% ======================================================================
\begin{thebibliography}{99}

\bibitem{Bellissard1992}
J.~Bellissard, B.~Iochum, E.~Scoppola, and D.~Testard,
\textit{Spectral properties of one-dimensional quasi-crystals},
Commun.\ Math.\ Phys.\ \textbf{125} (1989), 527--543.

\bibitem{Brent1973}
R.~P.~Brent,
\textit{Algorithms for Minimization without Derivatives},
Prentice-Hall, Englewood Cliffs, NJ, 1973.

\bibitem{Damanik2017}
D.~Damanik,
\textit{Schr\"{o}dinger operators with dynamically defined potentials},
Ergodic Theory Dynam.\ Systems \textbf{37} (2017), 1681--1764.

\bibitem{Duderstadt1976}
J.~J.~Duderstadt and L.~J.~Hamilton,
\textit{Nuclear Reactor Analysis},
Wiley, New York, 1976.

\bibitem{Grossi2026DNLS}
P.~Nogueira Grossi,
\textit{Differential Nonlinear Robustness of Critical States in
Fibonacci and Tribonacci Substitution Chains},
G6 LLC (2026), Zenodo.
\href{https://doi.org/10.5281/zenodo.20032853}{doi:10.5281/zenodo.20032853}.

\bibitem{GrossiRepo2026}
P.~Nogueira Grossi,
\textit{Atratores: numerical code repository},
GitHub (2026).
\href{https://github.com/grossi-ops/Atratores/tree/main/DNLS}{github.com/grossi-ops/Atratores}.

\bibitem{Krebbekx2023}
J.~P.~J.~Krebbekx, A.~Moustaj, K.~Dajani, and C.~de~Morais~Smith,
\textit{Multifractal properties of tribonacci chains},
Phys.\ Rev.\ B \textbf{108} (2023), 104204.
\href{https://doi.org/10.1103/PhysRevB.108.104204}{doi:10.1103/PhysRevB.108.104204}.

\bibitem{Allaire1999}
G.~Allaire and G.~Bal,
\textit{Homogenization of the criticality spectral equation in neutron transport},
ESAIM: Math.\ Model.\ Numer.\ Anal.\ \textbf{33} (1999), 721--746.

\bibitem{Varma2019}
V.~K.~Varma, C.~de~Tomasi, and M.~M\"{u}ller,
\textit{Diffusion in quasiperiodic Fibonacci chains},
Phys.\ Rev.\ B \textbf{100} (2019), 085105.

\bibitem{Shechtman1984}
D.~Shechtman, I.~Blech, D.~Gratias, and J.~W.~Cahn,
\textit{Metallic phase with long-range orientational order and no
translational symmetry},
Phys.\ Rev.\ Lett.\ \textbf{53} (1984), 1951--1953.

\bibitem{Squires1978}
G.~L.~Squires,
\textit{Introduction to the Theory of Thermal Neutron Scattering},
Cambridge University Press, Cambridge, 1978.

\bibitem{Stacey2007}
W.~M.~Stacey,
\textit{Nuclear Reactor Physics},
2nd ed., Wiley-VCH, Weinheim, 2007.

\end{thebibliography}

\end{document}
